
\chapter{QR分解}
\begin{definition}[正交单位向量组]
向量组 $\bm q_1,\ldots,\bm q_n$ 是\emph{正交单位}，若
$$
\bm q_i^{\mathsf T}\bm q_j=
\begin{cases}
0,& i\ne j \quad\text{（正交向量）},\\[2pt]
1,& i=j \quad\text{（单位向量：}\lVert \bm q_i\rVert=1\text{）}.
\end{cases}
$$
一个有正交列的矩阵都将被记为特征字母 $Q$。
\end{definition}

矩阵 $Q$ 很容易处理，因为
$$
Q^{\mathsf T}Q=I .
$$
在矩阵语言中再一次说明：若列向量 $\bm q_1,\ldots,\bm q_n$ 正交单位，
则 $Q$ 的列正交单位；注意 $Q$ 不一定需要是方阵（列正交即可，称为半正交矩阵\\矩阵 $Q$ 是方阵时，若满足
\[
Q^{\mathsf T} Q = I ,
\]
则有 $Q^{\mathsf T} = Q^{-1}$，也就是说转置矩阵等于逆矩阵。

进一步说明：当 $Q$ 是矩形矩阵时，只要 $Q^{\mathsf T} Q = I$，此时 $Q^{\mathsf T}$ 只是 $Q$ 的\emph{左逆矩阵}。若 $Q$ 是方阵，我们还可以得到
\[
Q Q^{\mathsf T} = I ,
\]
因此 $Q^{\mathsf T}$ 既是左逆也是右逆，即为 $Q$ 的\emph{双边逆矩阵}。在这种情形下，$Q$ 的行向量也像列向量一样正交归一，我们称这样的方阵 $Q$ 为\emph{正交矩阵}。\\
\begin{dxtips}[正交化的好处需要牢记]只有弄清楚了我们分解的目的才能聚焦在分解的手段
    
\end{dxtips}
一个 置换矩阵 是一个方阵，它的每一行、每一列里都有 且仅有一个元素是 1，其余元素全是 0。

换句话说，置换矩阵是单位矩阵 I 的行或列经过某种 重新排列（permutation） 得到的矩阵。
置换矩阵都是正交矩阵
\begin{example}[反射]
若 $u$ 是任意的单位向量，令
\[
Q = I - 2uu^{\mathsf T}.
\]
注意 $uu^{\mathsf T}$ 是一个矩阵，而 $u^{\mathsf T}u$ 是一个数字。由于 $\lVert u \rVert^2 = 1$，可得

\[
Q^{\mathsf T} = I - 2uu^{\mathsf T} = Q, \qquad 
Q^{\mathsf T}Q = I - 4uu^{\mathsf T} + 4uu^{\mathsf T}uu^{\mathsf T} = I. \tag{2}
\]

因此，反射矩阵 $I-2uu^{\mathsf T}$ 对称而且正交。如果把它平方，你会得到单位矩阵：
\[
Q^2 = Q^{\mathsf T}Q = I.
\]

这意味着通过镜子反射两次会回到原始状态，就好像 $(-1)^2=1$ 一样。注意在式 (2) 的 $4uu^{\mathsf T}uu^{\mathsf T}$ 中，利用了 $u^{\mathsf T}u=1$。
\end{example}   
\begin{dxtips}
    正交矩阵还能避免求逆矩阵的麻烦。
\end{dxtips}
旋转会保有每个向量的长度，反射(reflections)也是，置换(permutations)也是，用任何正交矩阵去乘也一样——长度与角度没有改变。  

证明：  
\[
\lVert Qx \rVert^2 \;=\; \lVert x \rVert^2 ,
\]
因为
\[
(Qx)^{\mathsf T}(Qx) \;=\; x^{\mathsf T} Q^{\mathsf T} Q x \;=\; x^{\mathsf T} I x \;=\; x^{\mathsf T}x .
\]

\begin{theorem}[正交矩阵保持长度与内积]
若 $Q$ 有正交单位列（即 $Q^{\mathsf T}Q = I$），它会保持向量的长度不变：  
\[
\text{对于每个向量 } x, \quad \lVert Qx \rVert = \lVert x \rVert. \tag{3}
\]

同时，$Q$ 也会保持点积：
\[
(Qx)^{\mathsf T}(Qy) \;=\; x^{\mathsf T} Q^{\mathsf T} Q y \;=\; x^{\mathsf T} y .
\]

这只应用了 $Q^{\mathsf T}Q = I$！
\end{theorem}
\section{转置子空间等知识回顾}
\begin{theorem}[微积分与转置]
    让我插入一点点微积分，这是极其重要的，否则我不会在讲线性代数时停下来。
（实际上，这里是针对函数 $x(t)$ 的线性代数。）
矩阵在这里变成了一个导数运算：$A = \dfrac{d}{dt}$。
要找到这个特殊 $A$ 的转置（更准确地说是伴随算子 adjoint），
我们需要先定义函数 $x(t)$ 与 $y(t)$ 的内积。

\[
x^{\mathsf T}y = (x,y) = \int_{-\infty}^{\infty} x(t) y(t)\, dt.
\]

因此，内积从有限维的求和 $\sum x_k y_k$ 变成了函数的积分。

---

矩阵的转置满足
\[
(Ax,y) = (x,A^{\mathsf T}y).
\]

对于 $A = \dfrac{d}{dt}$，我们有
\[
(Ax,y) = \int_{-\infty}^{\infty} \frac{dx}{dt}\, y(t)\, dt
= \int_{-\infty}^{\infty} x(t)\,\Bigl(-\frac{dy}{dt}\Bigr)\, dt
= (x, A^{\mathsf T}y).
\]

这里用了分部积分法（integration by parts）。  
在这个过程中，导数从第一个函数 $x(t)$ 转移到第二个函数 $y(t)$，并且会产生一个负号。  
这告诉我们：**导数的转置是它的相反数**。

---

因此，导数算子是反对称的（antisymmetric）：
\[
A = \frac{d}{dt}, \qquad A^{\mathsf T} = -\frac{d}{dt}.
\]

对称矩阵满足 $S^{\mathsf T} = S$，而反对称矩阵满足 $A^{\mathsf T} = -A$。  
线性代数中也包含微分与积分（见第 8 章），因为它们都是线性算子。
\end{theorem}
套用分部积分公式：
\[
\int u'(t)\,v(t)\,dt
= \Big[u(t)\,v(t)\Big]_{-\infty}^{\infty} - \int u(t)\,v'(t)\,dt.
\]

这里令
\[
u(t) = x(t), \quad u'(t) = \frac{dx}{dt}, \quad v(t) = y(t).
\]

那么：
\[
\int \frac{dx}{dt}\, y(t)\,dt
= \Big[ x(t)\,y(t)\Big]_{-\infty}^{\infty} - \int x(t)\,\frac{dy}{dt}\,dt.
\]

\medskip

如果假设边界项 $\Big[x(t)y(t)\Big]_{-\infty}^\infty=0$ 
（比如 $x(t),y(t)$ 在无穷远处衰减到 0，或者它们满足合适的边界条件），
那么就得到：
\[
\int_{-\infty}^\infty \frac{dx}{dt}\,y(t)\,dt
= -\int_{-\infty}^\infty x(t)\,\frac{dy}{dt}\,dt.
\]

\medskip

这就是书里写的：
\[
\int_{-\infty}^\infty \frac{dx}{dt}\,y(t)\,dt
= \int_{-\infty}^\infty x(t)\,\Bigl(-\frac{dy}{dt}\Bigr)\,dt.
\]

它说明了：导数算子的伴随（转置）就是带一个负号的导数算子。\begin{proposition}
    微分的反对称也可以应用到中心差分矩阵：
\[
A =
\begin{bmatrix}
0 & 1 & 0 & 0 \\
-1 & 0 & 1 & 0 \\
0 & -1 & 0 & 1 \\
0 & 0 & -1 & 0
\end{bmatrix}
\quad
\text{转置成}\quad
A^{\mathsf T} =
\begin{bmatrix}
0 & -1 & 0 & 0 \\
1 & 0 & -1 & 0 \\
0 & 1 & 0 & -1 \\
0 & 0 & 1 & 0
\end{bmatrix}
= -A
\]

前向差分矩阵的转置会变成反向差分矩阵，相当于乘 $-1$。  
在微分方程式中，第二导数（加速度）是对称的，第一导数（正比于速度的阻尼）是反对称的。
\end{proposition}
$A^{\mathsf T}A$ 的含义
1.~\textbf{代数性质}  
\begin{itemize}
  \item 对称性：$(A^{\mathsf T}A)^{\mathsf T} = A^{\mathsf T}A$；
  \item 半正定性：
  \[
  x^{\mathsf T}(A^{\mathsf T}A)x = (Ax)^{\mathsf T}(Ax) = \|Ax\|^2 \geq 0 ,
  \]
  因此特征值均为非负。
\end{itemize}

\medskip

2.~\textbf{几何意义}  
$A: \mathbb R^n \to \mathbb R^m$，而 $A^{\mathsf T}A: \mathbb R^n \to \mathbb R^n$。  
它的作用可以理解为：先通过 $A$ 把向量映射到输出空间，再通过 $A^{\mathsf T}$ 拉回输入空间。  
因此 $A^{\mathsf T}A$ 定义了一种新的“长度度量”：
\[
\|x\|_{A}^2 := x^{\mathsf T}(A^{\mathsf T}A)x = \|Ax\|^2 .
\]

\medskip

3.~\textbf{最小二乘中的作用}  
正规方程
\[
A^{\mathsf T}A \hat{x} = A^{\mathsf T}b
\]
说明 $A^{\mathsf T}A$ 扮演了 Gram 矩阵的角色，它把“残差最小化”的几何条件转化为代数方程。

\medskip

4.~\textbf{数据分析中的意义}  
$A^{\mathsf T}A$ 又称 Gram 矩阵或协方差矩阵（标准化后）。  
它记录了列向量之间的内积关系：
\[
(A^{\mathsf T}A)_{ij} = a_i^{\mathsf T}a_j,
\]
并且是 PCA（主成分分析）等方法的核心。

\medskip

\textbf{总结：} $A^{\mathsf T}A$ 既是对称半正定的 Gram 矩阵，反映了列向量的相关性，又是最小二乘与 PCA 的核心工具。
\begin{theorem}[转置与内积的意义]
    \[
x^{\mathsf T}y \;\; \text{是一个数字}, 
\qquad 
xy^{\mathsf T} \;\; \text{是一个矩阵}.
\]

在量子力学中常写作 $\langle x|y \rangle$ （内积）与 $|x\rangle \langle y|$ （外积）。

\medskip
宇宙可能是被线性代数管理的，下面三个例子对于内积有不同的意义：

\[
\begin{array}{lll}
\text{来自力学} & \quad \text{功 (work)} = (\text{位移})(\text{力}) &= x^{\mathsf T} f, \\[6pt]
\text{来自电路} & \quad \text{热损失} = (\text{压降})(\text{电流}) &= e^{\mathsf T} y, \\[6pt]
\text{来自经济} & \quad \text{收入} = (\text{数量})(\text{价格}) &= q^{\mathsf T} p.
\end{array}
\]
\end{theorem}
\begin{definition}[列空间、行空间与零空间]
还要再看
设 $A \in \mathbb{R}^{m\times n}$。

\begin{itemize}
  \item \textbf{列空间（Column Space）}  
  \[
  \mathcal{C}(A) = \{\, A x : x \in \mathbb{R}^n \,\} \subseteq \mathbb{R}^m.
  \]  
  它由 $A$ 的列向量张成，表示所有可能的输出 $Ax$。  
  几何意义：$A$ 把 $\mathbb{R}^n$ 中的向量映射到 $\mathbb{R}^m$，结果必定落在列空间中。  

  \item \textbf{行空间（Row Space）}  
  \[
  \mathcal{C}(A^{\mathsf T}) = \{\, A^{\mathsf T} y : y \in \mathbb{R}^m \,\} \subseteq \mathbb{R}^n.
  \]  
  它由 $A$ 的行向量张成，等价于 $A^{\mathsf T}$ 的列空间。  
  几何意义：行空间包含了 $A$ 对输入向量 $x$ 所进行的“线性约束”的所有可能组合。 也就是Ax=b里面每一行x的系数限制x之间的关系。 

  \item \textbf{零空间（Null Space）}  
  \[
  \mathcal{N}(A) = \{\, x \in \mathbb{R}^n : A x = 0 \,\}.
  \]  
  它是所有被 $A$ 映射到零向量的输入向量的集合。  
  几何意义：零空间表示了“失效方向”——这些方向上的输入完全被 $A$ 消去。
\end{itemize}
\end{definition}

\begin{proposition}[三者对比]
\begin{itemize}
  \item 列空间 $\mathcal{C}(A) \subseteq \mathbb{R}^m$，行空间 $\mathcal{C}(A^{\mathsf T}) \subseteq \mathbb{R}^n$，零空间 $\mathcal{N}(A) \subseteq \mathbb{R}^n$。
  \item 维数关系：
  \[
  \dim \mathcal{C}(A) = \dim \mathcal{C}(A^{\mathsf T}) = \mathrm{rank}(A).
  \]
  设 $A \in \mathbb{R}^{m\times n}$，那么
\[
\dim \mathcal{C}(A) + \dim \mathcal{N}(A) = n.
\]
  \item 正交关系：  
  零空间 $\mathcal{N}(A)$ 与行空间 $\mathcal{C}(A^{\mathsf T})$ 正交；  
  零空间 $\mathcal{N}(A^{\mathsf T})$ 与列空间 $\mathcal{C}(A)$ 正交。
\end{itemize}
\end{proposition}





\section{最小二乘法}
$Ax=b$ 经常是无解，一般的理由是方程式太多。矩阵 $A$ 的行比列多，方程式的个数大于未知数的个数（$m>n$）。
这 $n$ 个列只生成了 $m$ 维空间的一小部分，除非所有的量测值都很完美，否则 $b$ 会落在 $A$ 的列空间之外。
消元法得到不可能存在的方程式然后停止，但我们不能停止，因为量测值必然包含杂讯。

\textbf{重复：} 我们不可能一直把误差 $e=b-Ax$ 降为零；当 $e=0$ 时，$x$ 是 $Ax=b$ 的\emph{确切解}。
当 $e$ 的长度尽可能地小，$\hat{x}$ 就是“\emph{最小二乘解}”（least squares solution）。
本段的目标是计算并使用 $\hat{x}$，许多实际问题都需要它。

前一段强调投影 $p$，本段强调最小二乘解 $\hat{x}$；二者通过 $p=A\hat{x}$ 相关联。
基本方程式仍然是正规方程：$A^{\mathsf T} A x = A^{\mathsf T} b$。
下面给出一个非正式的得到“正规方程式”的办法：

\begin{note}
当 $Ax=b$ 无解时，左乘 $A^{\mathsf T}$ 然后求解
\[
A^{\mathsf T}A\,x \;=\; A^{\mathsf T}b .
\]
\end{note}
我们希望用一条直线
\[
y \approx ax + b
\]
来拟合一组点。将参数写成向量
\[
\theta = 
\begin{bmatrix}
a \\[4pt] b
\end{bmatrix},
\]
并将数据点 $(x_i,y_i)$ 写成矩阵形式：
\[
A =
\begin{bmatrix}
x_1 & 1 \\
x_2 & 1 \\
\vdots & \vdots \\
x_m & 1
\end{bmatrix},
\qquad
y =
\begin{bmatrix}
y_1 \\ y_2 \\ \vdots \\ y_m
\end{bmatrix}.
\]

于是所有观测点可以统一写作
\[
A \theta \;\approx\; y ,
\]
即
\[
\begin{bmatrix}
x_1 & 1 \\
x_2 & 1 \\
\vdots & \vdots \\
x_m & 1
\end{bmatrix}
\begin{bmatrix}
a \\[4pt] b
\end{bmatrix}
=
\begin{bmatrix}
ax_1+b \\ ax_2+b \\ \vdots \\ ax_m+b
\end{bmatrix}
\;\approx\;
\begin{bmatrix}
y_1 \\ y_2 \\ \vdots \\ y_m
\end{bmatrix}.
\]

由于一般情况下方程组无解，我们改为最小二乘问题：
\[
\hat{\theta} = \arg\min_{\theta} \|A\theta - y\|^2 ,
\]
其解为正规方程
\[
\hat{\theta} = (A^{\mathsf T}A)^{-1} A^{\mathsf T} y .
\]
\begin{dxtips}[ 数值分析课上关于最小二乘法的原理]
    这是经典的最小二乘误差函数：
\[
E(a_0,a_1) = \sum_{i=1}^n \big( y_i - a_0 - a_1 x_i \big)^2 .
\]

用导数方法求最小值（原理关键！

我们令误差函数对 $a_0, a_1$ 的偏导数为 0，找极小值点：
\[
\frac{\partial E}{\partial a_0} = 0, \qquad \frac{\partial E}{\partial a_1} = 0.
\]

这是一个标准的“求极值”步骤。展开以后可以得到两个线性方程：
\[
\begin{cases}
n a_0 + a_1 \sum x_i = \sum y_i, \\[6pt]
a_0 \sum x_i + a_1 \sum x_i^2 = \sum x_i y_i.
\end{cases}
\]
\end{dxtips}
直线拟合问题写成矩阵形式：
\[
A\theta \approx y, \quad 
A=
\begin{bmatrix}
x_1 & 1 \\
x_2 & 1 \\
\vdots & \vdots \\
x_m & 1
\end{bmatrix},\quad
\theta=\begin{bmatrix}a\\b\end{bmatrix},\quad
y=\begin{bmatrix}y_1\\y_2\\\vdots\\y_m\end{bmatrix}.
\]

因为 $m>2$ 一般无解，改为最小二乘，得到正规方程
\[
A^{\mathsf T}A\theta=A^{\mathsf T}y,
\]
即
\[
\begin{bmatrix}
\sum x_i^2 & \sum x_i \\
\sum x_i   & m
\end{bmatrix}
\begin{bmatrix}
a \\ b
\end{bmatrix}
=
\begin{bmatrix}
\sum x_i y_i \\
\sum y_i
\end{bmatrix}.
\]

于是得到两条方程：
\[
\begin{cases}
a\sum x_i^2 + b\sum x_i = \sum x_i y_i, \\[6pt]
a\sum x_i + b m = \sum y_i ,
\end{cases}
\]
其解 $(a,b)$ 正是使残差平方和最小的拟合直线参数。
\begin{proposition}[从矩阵理解和从导数理解的关系]
    我们要最小化的目标函数是
\[
E(\theta) = \|A\theta - y\|^2 = (A\theta - y)^{\mathsf T}(A\theta - y),
\]
这是最小二乘的定义。

\medskip

\textbf{2. 求偏导 = 0}

对 $\theta$ 求梯度：
\[
\nabla_\theta E = 2A^{\mathsf T}(A\theta - y).
\]

设导数 $=0$，则
\[
A^{\mathsf T}(A\theta - y) = 0,
\]
这就得到正规方程
\[
A^{\mathsf T}A \theta = A^{\mathsf T} y.
\]

因此，“偏导数 $=0$” 推出的条件，和 “左乘 $A^{\mathsf T}$ 保证残差正交”其实是同一个式子。
\end{proposition}
\begin{proposition}[如何从投影理解优化目标]
矩阵 $A$ 的所有列向量可以生成一个列空间，向量 $Ax$ 一定在这个列空间内，但向量 $b$ 可能不在列空间内。
$Ax$ 与 $b$ 的差就是一个新的向量 $r$，称为残差，可以用平行四边形法则理解。
为了使 $r$ 的长度最短，必须要求它垂直于由 $A$ 的列向量生成的空间。
如果满足这个条件，就有
\[
A^{\mathsf T} r = 0,
\]
此时残差的长度也达到最小。
几何上可以想象为：空间外一点到平面的最短距离就是垂直距离，而把 $b$ 投影到由 $A$ 的列向量张成的平面上后，就存在相应的 $x$，使得 $Ax$ 恰好等于这个投影向量。
\end{proposition}
\begin{dxtips}
    从几何，代数和微积分都能理解这个方程的目的。没多一种维度你对问题的理解就更加深刻。
\end{dxtips}
\begin{dxtips}
    

矩阵乘法的意义     \\
✅ 一句话：谁在左边，谁就“主导”最后结果落在哪个空间；右边的对象只决定“怎么组合”，因此，当我们同时左乘 $A^{\mathsf T}$ 时，等式的两边都被映射到了 
$A^{\mathsf T}$ 所作用的空间中。选择左乘 $A^{\mathsf T}$ 不是随意的，而是因为它唯一能把几何条件$r \perp \mathcal C(A)$（残差 垂直列空间）翻译成代数形式。
\end{dxtips}

\\
列空间与行空间


\[
(Ax) \cdot y = x \cdot (A^{\mathsf T} y)
\]

$A$ 生成的空间是“输出空间里的列空间”——列向量为基生成的空间，$A^{\mathsf T}$ 生成的空间是“输入空间里的行空间”——行向量生成空间。它们的维数相等，但位置不同，通过点积保持联系。
所以，输入空间就是 $\mathbb{R}^n$（因为你只能送 n 维向量进去）也就是乘的x必须是一个n维列向量，
输出空间就是$ \mathbb{R}^m$（因为出来的永远是 m 维）。

\section{QR分解操作}
QR 分解的核心目的是将矩阵 $A$ 分解为
\[
A = QR,
\]
其中 $Q$ 是列向量正交归一的矩阵，$R$ 是上三角矩阵。换句话说，QR 分解就是把原来可能不正交的基向量转化为一组正交单位向量，并将变化过程中的比例关系记录在 $R$ 中。和前面的矩阵乘法概念类似，能清晰的看出R对原本的正交基有怎么样的拉伸和旋转。

\begin{proposition}
当 $Q$ 是方阵且 $m=n$ 时，子空间就是整个空间。此时有
\[
Q^{\mathsf T} = Q^{-1},
\]
并且
\[
\hat{x} = Q^{\mathsf T} b \quad \text{与} \quad x = Q^{-1} b
\]
完全相同，因此这个解是确切的。

此外，向量 $b$ 在整个空间上的投影就是它自身：
\[
p = b,
\]
并且投影矩阵满足
\[
P = Q^{\mathsf T} Q = I.
\]
这样我们从原来的两边同乘$A^{\mathsf T}$ 变成两边同乘$Q^{\mathsf T}$
	乘 $A^{\mathsf T}A$ 的方法会 放大条件数(奇异值变成平方)，误差更大。
	•	用 $Q^{\mathsf T}$ 投影不会改变长度或角度（正交变换保持数值稳定）。R作为上三角矩阵天然好解方程。
\end{proposition}
\begin{proposition}[正交矩阵与傅立叶变换]
    当 $p=b$，我们的公式把 $b$ 从它的一维投影中聚集起来。
若 $q_1,\dots,q_n$ 是整个空间的一组正交单位基底，则 $Q$ 是方形，
每个 $b=QQ^{\mathsf T}b$ 是它沿着 $q$'s 的分量的总和：

\[
b = q_1(q_1^{\mathsf T} b) + \cdots + q_n(q_n^{\mathsf T} b).
\]

\noindent
\textbf{转换} $QQ^{\mathsf T}=I$ 是傅里叶级数的基础，也是所有应用数学中伟大“转换”的基础。
他们把 $b$ 或函数 $f(x)$ 打碎成垂直的片段，借由上式中的片段，逆转换把 $b$ 与 $f(x)$ 再次一起恢复。
\end{proposition}
\begin{dxtips}[施密特正交化]
    看见转置要想到点积，看见点积要想到投影。  
    如何把三个向量 $a,b,c$ 正交化呢？  
    方法就是：把每一个向量在其他向量方向上的分量减掉。  

    \medskip
    格莱姆–施密特的第一步：
    \[
        B \;=\; b \;-\; \frac{A^{\mathsf T} b}{A^{\mathsf T} A} \, A,
        \tag{7}
    \]

    其中投影系数为
    \[
        \frac{\langle b, A \rangle}{\langle A, A \rangle}
        \;=\;
        \frac{A^{\mathsf T} b}{A^{\mathsf T} A}.
    \]
    两边左乘A的转置即可验证A与B的正交。c同理去掉a和B的分量产生C。最后a，B，C再除以自己的模长即可。\textbf{一句话总结}

Gram-Schmidt 正交化 $=$ 把 $A$ 的列向量逐个「扣掉已有正交方向的分量」。这个扣除的系数恰好填入 $R$ 的上三角位置，剩下的单位化结果填入 $Q$ 的列。$A = QR$ 不是额外的结论，而是 Gram-Schmidt 过程的直接重写。
\end{dxtips}
\begin{proposition}[格莱姆–施密特的关键步骤]
第一步是
\[
q_1 = \frac{a}{\|a\|},
\]
（其他向量不参与）。第二步是方程式 (7)，其中 $b$ 是 $A$ 与 $B$ 的组合，在这一步骤 $c$ 与 $q_3$ 没有参与。  
下一个向量不参与正是格莱姆–施密特的关键点：这就是为什么下三角可以直接写0。

\begin{itemize}
    \item 向量 $a$ 与 $A$ 与 $q_1$ 沿着同一条直线；
    \item 向量 $a,b$ 与 $A,B$ 与 $q_1,q_2$ 在同一个平面；
    \item 向量 $a,b,c$ 与 $A,B,C$ 与 $q_1,q_2,q_3$ 在同一个子空间（维度 3）。
\end{itemize}

在每一步，$a_1,\dots,a_k$ 是 $q_1,\dots,q_k$ 的组合，后面的 $q$ 不参与。  
关联矩阵 $R$ 是三角形，我们有
\[
A = QR.
\]

\[
\bigl[\, a \;\; b \;\; c \,\bigr]
= 
\bigl[\, q_1 \;\; q_2 \;\; q_3 \,\bigr]
\begin{bmatrix}
q_1^{\mathsf T} a & q_1^{\mathsf T} b & q_1^{\mathsf T} c \\[6pt]
0 & q_2^{\mathsf T} b & q_2^{\mathsf T} c \\[6pt]
0 & 0 & q_3^{\mathsf T} c
\end{bmatrix},
\qquad
\text{或简记为 } \; A = QR.
\tag{9}
\]
\end{proposition}
 \begin{note}
设 $A$ 的列向量是 $a_1, a_2, a_3$，做 Gram-Schmidt 正交化：
\[
q_1 = \frac{a_1}{\|a_1\|}, \quad
q_2 = \frac{a_2 - (a_2 \cdot q_1)q_1}{\|\cdots\|}, \quad
q_3 = \frac{a_3 - (a_3 \cdot q_1)q_1 - (a_3 \cdot q_2)q_2}{\|\cdots\|}
\]
\textbf{反过来}用 $q_i$ 表示原列向量：
\[
a_1 = \|a_1\| q_1, \quad
a_2 = (a_2 \cdot q_1)q_1 + \|\cdots\| q_2, \quad
a_3 = (a_3 \cdot q_1)q_1 + (a_3 \cdot q_2)q_2 + \|\cdots\| q_3
\]
写成矩阵形式即 $A = QR$：
\[
\begin{bmatrix} | & | & | \\ a_1 & a_2 & a_3 \\ | & | & | \end{bmatrix}
=
\begin{bmatrix} | & | & | \\ q_1 & q_2 & q_3 \\ | & | & | \end{bmatrix}
\begin{bmatrix} r_{11} & r_{12} & r_{13} \\ 0 & r_{22} & r_{23} \\ 0 & 0 & r_{33} \end{bmatrix}
\]
\textbf{$R$ 为何上三角？}\ 因为 $a_k$ 正交化时只涉及前 $k$ 个基向量，第 $k$ 列仅 $r_{1k},\ldots,r_{kk}$ 非零，天然产生上三角结构。
\end{note}
\begin{theorem}[格莱姆–施密特正交化]
从无关向量 $a_1, \dots, a_n$ 出发，格莱姆–施密特方法能够构造出正交单位向量 
$q_1, \dots, q_n$。这些列向量构成的矩阵 $Q$ 满足
\[
A = QR,
\]
其中
\[
R = Q^{\mathsf T} A
\]
是一个上三角矩阵，因为后面的 $q$ 与较早的 $a$ 的内积为零（正交性）。
\end{theorem}
\begin{equation}
\text{最小二乘:} \quad 
R^{\mathsf T} R \hat{x} = R^{\mathsf T} Q^{\mathsf T} b
\quad \text{或} \quad
R \hat{x} = Q^{\mathsf T} b
\quad \text{或} \quad
\hat{x} = R^{-1} Q^{\mathsf T} b
\tag{10}
\end{equation}
\section{QR分解迭代法-数值分析}
基本的步骤是分解 $A$ 成为 $QR$，其中 $A$ 的特征值是我们想要的。回忆格莱姆–施密特 (段落 4.4)，$Q$ 有正交单位列与 $R$ 是三角形。特征值的关键概念是：反转 $Q$ 与 $R$。新的矩阵 (相同 $\lambda$’s) 是 $A_1 = RQ$，因为 $A = QR$ 与 $A_1 = Q^{-1}AQ$ 相似，所以 $RQ$ 的特征值没有改变：
\[
A_1 = RQ \ \text{有相同的特征值 } \lambda, 
\quad QRx = \lambda x 
\ \Rightarrow\ 
RQ(Q^{-1}x) = \lambda (Q^{-1}x).
\tag{12}
\]

继续这个程序。把新矩阵 $A_1$ 分解成 $Q_1R_1$，然后反转因子变成 $R_1Q_1$，这是相似矩阵 $A_2$ 而且再次没有改变特征值。惊人的是这些特征值开始出现在对角线，很快的 $A_4$ 的最后单元得到一个精确的特征值。这个情况下，我们移除最后一行与最后一列得到一个较小的矩阵，然后继续寻找下一个特征值。


这个现象的原因是
\[
A_k = (Q_1 Q_2 \cdots Q_k)^{\!T}\, A_0\, (Q_1 Q_2 \cdots Q_k).
\]

记累计的正交变换
\[
Q^{(k)} := Q_1 Q_2 \cdots Q_{k-1} \quad (\text{正交}),
\]

于是
\[
A_k = (Q^{(k)})^{\!T}\, A_0\, Q^{(k)}.
\]

这说明：幂迭代是找的一个特征向量，QR 迭代是找一组特征向量做成的正交基
\\
\begin{proposition}
    两个额外的概念使得这个方法成功，一个是在分解成 $QR$ 之前，把矩阵平移 $I$ 的倍数，
然后 $RQ$ 再平移回去得到 $A_{k+1}$：

\[
A_k - c_k I = Q_k R_k, 
\qquad 
A_{k+1} = R_k Q_k + c_k I.
\]

于是 $A_{k+1}$ 与 $A_k$ 有相同的特征值，而且与原始的 $A_0=A$ 有相同的特征值。
一个优质的平移是选择 $c$ 靠近一个（未知的）特征值。
这个更精确的特征值出现在 $A_{k+1}$ 的对角线上——这告诉我们下一个到 $A_{k+2}$ 的步骤应该选择的较好 $c$。

---

第二个概念是在 QR 方法开始之前先获得非对角线的零，一个消元步骤 $E$ 会做这样的事情，
或者是一个吉文斯旋转（Givens rotation）。
但是不要忘了 $E^{-1}$（否则 $\lambda$ 会改变）：

\[
E A E^{-1}
=
\begin{bmatrix}
1 & 1 & 2 \\
0 & 1 & 4 \\
-1 & 1 & 6
\end{bmatrix}
\begin{bmatrix}
1 & 2 & 3 \\
1 & 4 & 5 \\
1 & 6 & 7
\end{bmatrix}
\begin{bmatrix}
1 & 1 & 1 \\
0 & 1 & 1 \\
0 & 0 & 1
\end{bmatrix}
=
\begin{bmatrix}
1 & 5 & 3 \\
1 & 9 & 5 \\
0 & 4 & 2
\end{bmatrix}
\quad \text{（相同的 $\lambda$'s）}.
\]

我们必须在 \textbf{次对角线} 留下非零数 $1$ 与 $4$，更多的 $E$ 可以移除它们，
但是 $E^{-1}$ 会再次填满。这是一个 \textbf{海森堡矩阵}（一个非零的次对角线）。

---

在 QR 方法过程中，左下角的零仍然维持是零。
因此，每个 QR 分解的运算次数可以从 $O(n^3)$ 降到 $O(n^2)$。
\end{proposition}
\begin{proposition}[幂迭代与逆迭代]
最大特征值对应的向量会占据主导地位因为同时提出$\lambda_1^{k}$ ，除了最大的一项，其他的$\lambda_n/\lambda_1$比值都小于1，k趋紧无穷的时候他们都趋紧于0.  
\end{proposition}
从任意向量 $u_0$ 开始。用 $A$ 乘它得到 $u_1$。再用 $A$ 乘，得到 $u_2$。  
如果 $u_0$ 是特征向量的线性组合，那么 $A$ 乘以每个特征向量 $x_i$ 时会放大 $\lambda_i$ 倍。  
经过 $k$ 步之后，我们得到 $(\lambda_i)^k$：

\[
u_k = A^k u_0 = c_1 (\lambda_1)^k x_1 + \cdots + c_n (\lambda_n)^k x_n. \tag{10}
\]

随着幂迭代的进行，最大的特征值会逐渐占据主导地位。向量 $u_k$ 会逐渐指向那个主特征向量 $x_1$。  
我们在马尔可夫矩阵中看到过这一点：

\[
A = \begin{bmatrix}
0.9 & 0.3 \\
0.1 & 0.7
\end{bmatrix}
\quad \text{其最大特征值 } \lambda_{\max} = 1, \text{ 对应的特征向量为 }
\begin{bmatrix} 0.75 \\ 0.25 \end{bmatrix}.
\]

从 $u_0$ 开始，并在每一步都用 $A$ 相乘：

\[
u_0 = \begin{bmatrix}1 \\ 0\end{bmatrix}, \quad
u_1 = \begin{bmatrix}0.9 \\ 0.1\end{bmatrix}, \quad
u_2 = \begin{bmatrix}0.84 \\ 0.16\end{bmatrix},
\]

逐渐趋向于

\[
u_\infty = \begin{bmatrix}0.75 \\ 0.25\end{bmatrix}.
\]

\medskip

收敛速度取决于第二大特征值 $\lambda_2$ 与最大特征值 $\lambda_1$ 的比值。  
我们并不希望 $\lambda_1$ 太小，而是希望 $\lambda_2/\lambda_1$ 较小。  
在这个例子中，$\lambda_2 = 0.6$，$\lambda_1 = 1$，因此收敛速度很好。  

对于大矩阵，往往会发生 $|\lambda_2/\lambda_1|$ 很接近 1 的情况，这时幂迭代就太慢了。

\medskip

有没有办法找到最小的特征值（这往往在应用中最重要）呢？  
答案是肯定的：用反幂迭代（inverse power method）：  
用 $A^{-1}$ 而不是 $A$ 去乘 $u_0$。  
由于我们从不真的去计算 $A^{-1}$，实际上我们解的是 $Au_1 = u_0$。  
通过保存 $LU$ 分解，下一步 $Au_2 = u_1$ 也可以快速完成。  
在第 $k$ 步时，有：

\[
Au_k = u_{k-1}.
\]
\[
u_k = A^{-k} u_0 = \frac{c_1 x_1}{(\lambda_1)^k} + \cdots + \frac{c_n x_n}{(\lambda_n)^k}. \tag{11}
\]

现在，最小的特征值 $\lambda_{\min}$ 起主导作用。  
当它非常小时，因子 $1/\lambda_{\min}^k$ 就会很大。  
为了加快速度，我们通过把矩阵平移为 $A - \lambda^* I$ 来使 $\lambda_{\min}$ 变得更小。  

这种平移不会改变特征向量。（$\lambda^*$ 可以取自 $A$ 的对角元，更好的选择是 Rayleigh 商 
\[
\frac{x^{\mathsf T} A x}{x^{\mathsf T} x}.
\]）

如果 $\lambda^*$ 接近 $\lambda_{\min}$，那么 $(A - \lambda^* I)^{-1}$ 就有一个非常大的特征值 $(\lambda_{\min} - \lambda^*)^{-1}$。  

每一步 \textit{平移反幂迭代} 都会把该特征向量乘以这个大数，从而使该特征向量迅速占据主导地位。
\chapter{LLL格基}

\begin{definition}[格与格基]
设 $b_1, b_2, \dots, b_n \in \mathbb{R}^m$ 是 $n$ 个线性无关向量。由它们生成的格定义为
\[
L(b_1, b_2, \dots, b_n) \;=\; \Bigl\{\, z_1 b_1 + z_2 b_2 + \cdots + z_n b_n \;\Big|\; z_i \in \mathbb{Z} \,\Bigr\}.
\]

称 $B = (b_1, b_2, \dots, b_n)$ 为格 $L(B)$ 的\emph{格基}。
注意，一个格可以有多组不同的基。
\end{definition}
\begin{remark}
与普通的向量空间基不同，格基规定系数必须是整数：
\[
L(B) = \{ z_1 b_1 + \cdots + z_n b_n \mid z_i \in \mathbb{Z} \}.
\]
因此，格是离散的点集，而不是连续的空间。
\end{remark}
\begin{proposition}[LLL 算法的核心思想]
LLL 算法不能追求完全正交（那样就不再是格基），
只能在“整数保持”和“尽量正交”之间折中：
\begin{itemize}
    \item \textbf{保证：} 所有变换都是整系数可逆矩阵（unimodular），因此格 
    \(\mathcal{L}(B)\) 保持不变；
    \item \textbf{目标：} 基向量不要太“偏斜”，做到“近似正交”，同时“尽量短”。
\end{itemize}
\end{proposition}


\subsection*{输入与目标}
设格基为
\[
B = (b_1, b_2, \dots, b_n) \in \mathbb{Z}^{m\times n}.
\]
目标：通过基变换得到 ``更短、更接近正交'' 的基。其中每一个向量里面全是整数，为最后的势函数(B)提供约束。
\[
B' = (b_1', b_2', \dots, b_n'),
\]
满足一定的正规化条件（LLL-reduced）。

\subsection*{Gram--Schmidt 正交化}
记 Gram--Schmidt 结果：
\[
b_i^* = b_i - \sum_{j=1}^{i-1} \mu_{j,i} b_j^*, 
\qquad 
\mu_{j,i} = \frac{\langle b_i, b_j^* \rangle}{\langle b_j^*, b_j^* \rangle}.
\]
于是 $b_i^*$ 是 $b_i$ 在 $\mathrm{span}\{b_1,\dots,b_{i-1}\}$ 的正交补上的投影。
\begin{dxtips}
 $\mu_{j,i}$是$b_i$中含有 $\mu_{j,i}$倍的$b_j$，如果这个系数很大，说明 $b_i$ 在 $b_j^*$ 方向上“粘连”得厉害 → 基向量会偏斜。
\end{dxtips}


\subsection*{步骤一：规模约化（Size Reduction）}
\[
b_i \;\leftarrow\; b_i - \lfloor \mu_{j,i} \rceil\, b_j, 
\quad \text{for } j<i,
\]
其中 $\lfloor \cdot \rceil$ 表示四舍五入到最近的整数。 
不断减去整数倍的前面向量，避免 ``过度粘连''，并且保证点仍在格内。让系数 $\mu_{j,i}$ 控制在
\[
|\mu_{j,i}| \leq \tfrac{1}{2}.
\]
这保证了基向量不会 ``偏斜''，得到更短正交。
\begin{definition}[Unimodular 矩阵]
一个整数矩阵 $U \in \mathbb{Z}^{n \times n}$ 若满足
\[
\det(U) = \pm 1,
\]
则称为 \emph{unimodular 矩阵}。
\end{definition}
\begin{proposition}[Unimodular 矩阵及其作用]
若 $U \in \mathbb{Z}^{n \times n}$ 且 $\det(U) = \pm 1$，则称 $U$ 为 \emph{unimodular 矩阵}。它有以下性质：

\begin{enumerate}
    \item \textbf{保持格不变.}  
    若 $B$ 是格基，则
    \[
    \mathcal{L}(B) = \{ Bz : z \in \mathbb{Z}^n \}, \qquad
    B' = BU,
    \]
    则有
    \[
    \mathcal{L}(B') = \mathcal{L}(B).
    \]
    因为 $U$ 可逆且整系数，保证 $\mathbb{Z}^n$ 在变换下仍映射回 $\mathbb{Z}^n$。

    \item \textbf{保持体积.}  
    平行多面体体积满足
    \[
    \det(BU) = \det(B)\det(U) = \pm \det(B),
    \]
    因此 $|\det(B)|$ 保持不变。

    \item \textbf{LLL 算法中的意义.}  
    在 LLL 算法中，所有基向量的调整都可表示为右乘 unimodular 矩阵：
    \begin{itemize}
        \item \emph{规模约化（Size Reduction）}：  
        \[
        b_i \;\leftarrow\; b_i - \lfloor \mu_{i,j} \rceil b_j, \quad j<i ,
        \]
        对应于右乘一个整系数初等矩阵，保证 $\mu_{i,j}$ 被控制在 $\bigl[-\tfrac{1}{2},\tfrac{1}{2}\bigr]$。
        \item \emph{交换步骤（Swap）}：  
        \[
        b_i \leftrightarrow b_{i-1},
        \]
        对应于交换两列，也就是 unimodular 矩阵。
    \end{itemize}
    这些操作保证始终在同一个格中进行换基，而不会改变格的本质。
\end{enumerate}

\noindent
因此，unimodular 矩阵就是“合法的换基操作”：它能改变基向量的形状，但不会改变格本身，也不会改变体积。
\end{proposition}

\subsection*{步骤二：Lovász 条件与交换}
在规模约化后，检查相邻两列是否满足 Lovász 条件：
\[
\|b_i^*\|^2 \;\ge\; (\delta - \mu_{i-1,i}^2)\,\|b_{i-1}^*\|^2,
\qquad \delta \in \Bigl(\tfrac{1}{4},1\Bigr),
\]
通常取 $\delta = \tfrac{3}{4}$。

\begin{itemize}
\item Lovasz条件大致保证当前工作的向量足够大成为我们的临时基向量\\在 Gram--Schmidt 正交化中，一个基向量 $b_i$ 会被投影到前面的正交补 $b_i^*$ 上。
\begin{itemize}
  \item 如果 $b_i^*$ 太短，说明 $b_i$ 在方向上几乎被前面基向量“解释掉了”；
  \item 那么整个基就很“塌”，缺乏独立性，导致表示格点的效率很差。
\end{itemize}

于是 LLL 算法要求：
\[
\|b_i^*\|^2 \;\;\text{不能太小，相对前一个正交向量 } b_{i-1}^* \text{ 而言要有下界。}
\]
  \item 若条件成立：继续处理 $i+1$ 列；
  \item 若条件不成立：交换 $b_{i-1}$ 与 $b_i$，并令 $i \gets \max(i-1,2)$，然后重新规模约化。
\end{itemize}
\begin{proposition}[LLL 算法的交换思想]
在施密特正交化中，按顺序处理基向量：前面的基先正交化，后面的基不断减去前面的投影。
LLL 的交换条件（Lovász 条件）就是在检查：

\begin{itemize}
    \item 前面的基向量是否 ``太差劲''（太长或太偏斜），
    \item 如果条件不满足，就交换 $b_k$ 和 $b_{k+1}$。
\end{itemize}

这样做的结果是：
\begin{itemize}
    \item 较短、方向更合适的向量被提前，作为参照基；
    \item 较差的向量被放到后面，被迫减去 ``好的'' 向量的分量。
\end{itemize}

\medskip
\textbf{Lovász 条件：}
\[
\delta \,\|b^*_k\|^2 \;\leq\; \|b^*_{k+1} + \mu_{k+1,k}\, b^*_k\|^2,
\]
其中 $b^*_i$ 是施密特正交向量，$\delta \in (1/4,1)$（通常取 $\delta=3/4$）。

\begin{itemize}
    \item 若条件成立，说明 $b_k$ ``够好''，无需交换；
    \item 若条件不成立，说明 $b_k$ 太差，就与 $b_{k+1}$ 交换。
\end{itemize}

\medskip
\textbf{类比：}  
可以把这一过程看作 ``选队长''：  
\begin{itemize}
    \item 好的队员（短而正交性好）先当队长；  
    \item 差的队员只能排在后面，被迫去配合前面更强的队员。  
\end{itemize}

这样，整个队伍（格基）就会更整齐、更短、更接近正交。
\end{proposition}
\begin{proposition}[LLL 与冒泡排序的类比]
在冒泡排序中，较大的元素（若从大到小排序）会一步步“冒”到前面：
\begin{itemize}
  \item 每次比较相邻两个，如果顺序错误，就交换。
  \item 经过多轮，最终最“大”的元素会被送到最前面。
\end{itemize}

在 LLL 算法中，同样存在类似的“冒泡”现象：
\begin{itemize}
  \item LLL 希望“更短、更正交”的向量排在前面。
  \item 每次检查 Lovász 条件
  \[
    \delta \,\|b_k^*\|^2 \;\leq\; \bigl\| b_{k+1}^* + \mu_{k+1,k} b_k^* \bigr\|^2,
  \]
  若不满足，就交换 $b_k \leftrightarrow b_{k+1}$。
  \item 交换后，那个“更好”的向量（更短、更正交的）会向前移动一步，
  就像冒泡排序里较大的元素逐步向前冒。
  \item 因此，经过不断交换，“最好的基向量”会逐渐被推到前面。
\end{itemize}
\end{proposition}
\begin{dxtips}[类比：LLL 的 Lovász 条件 vs ODE 的 Lipschitz 条件]
\[
\begin{array}{|c|c|}
\hline
\textbf{LLL 格基约化} & \textbf{常微分方程 (ODE)} \\ \hline
\text{Lovász 条件：保证相邻正交向量不能“太塌”} 
& \text{Lipschitz 条件：保证函数变化不能“太陡”} \\[6pt]
\delta \, \|b_k^*\|^2 \;\leq\; 
\|b_{k+1}^* + \mu_{k+1,k} b_k^*\|^2, \quad \delta \in (1/4,1) 
& |f(x,y_1) - f(x,y_2)| \;\leq\; L |y_1-y_2| \\[8pt]
\text{避免基向量几乎线性相关，保证格的稳定性} 
& \text{避免解不唯一，保证 ODE 解的稳定性} \\ \hline
\end{array}
\]
\end{dxtips}
\subsection*{步骤三：终止性与复杂度}
定义一个势函数
\[
\Phi(B) = \prod_{i=1}^n \|b_i^*\|^2,
\]
这个势函数就所有$b_i$模长的平方的乘积。它在每次交换后都会严格减小一个有界的整数量，因此算法在有限步后必然终止。

复杂度结果：LLL 在输入长度的多项式时间内完成（这是它在密码学应用中的关键优势）。
\begin{proposition}[LLL 势函数与格体积]
设 $B=[b_1,\dots,b_n]\in\mathbb{R}^{m\times n}$，Gram–Schmidt 正交化得
$b_1^*,\dots,b_n^*$。则
\[
\det(B^{\mathsf T}B) \;=\; \prod_{i=1}^n \|b_i^*\|^2 .
\]

因此定义的势函数
\[
\Phi(B) := \prod_{i=1}^n \|b_i^*\|^2
\]
正是格体积的平方。  

由于 LLL 算法中的变换均为 unimodular 矩阵作用（保持体积不变），
$\Phi(B)$ 始终与格的几何结构对应，并在交换或约化时严格下降，
保证算法收敛。
\end{proposition}
---



\begin{proposition}{有限步终止性的直观理解.}
在整格 $\mathcal{L}(B)\subseteq \mathbb{Z}^m$ 中，任意一组基 $B$ 生成的平行多面体
\[
P(B) = \Bigl\{\sum_{i=1}^n \alpha_i b_i : 0 \leq \alpha_i < 1 \Bigr\}
\]
的体积必为整数，因为它正好对应 $\mathbb{R}^m$ 中由单位立方格子堆叠出的“积木”个数。  
因此
\[
\det(\mathcal{L}(B)) = |\det(B)| \in \mathbb{Z}_{>0}.
\]

Gram–Schmidt 分解给出
\[
\prod_{i=1}^n \|b_i^*\| = |\det(B)|,
\]
于是势函数
\[
\Phi(B) = \prod_{i=1}^n \|b_i^*\|^2 \in \mathbb{Z}_{>0}.
\]

换句话说，每一个格基形成的平行多面体体积天然是整数，  
势函数就是它的“平方版”。每次交换都会严格减小 $\Phi(B)$，  
而 $\Phi(B)$ 只有有限多个整数值，因此 LLL 算法必在有限步后终止。
\end{proposition}
\begin{dxtips}
    离散的好处就是有限，至少相对连续来说好很多
\end{dxtips}

\subsection*{结果性质}
输出基 $(b_1',\dots,b_n')$ 满足：
\begin{itemize}
  \item \textbf{规模约化：} $\; |\mu_{j,i}| \le \tfrac12,\ \forall j<i$；
  \item \textbf{Lovász 条件：} $\; \|b_i^*\|^2 \ge (\delta-\mu_{i-1,i}^2)\|b_{i-1}^*\|^2$；
  \item \textbf{基向量较短且接近正交。}
\end{itemize}
\section{LLL最短向量的意义}
用 LLL 从数值 \(\alpha\) 识别整系数多项式：以 \(\alpha\approx\sqrt2\) 为例

{目标}
已知一个实数的近似值 \(\alpha\)（例如 \(\alpha=1.414213\)），
希望找到一个\emph{整系数多项式}
\[
f(x)=c_0+c_1x+\cdots+c_dx^d\in\mathbb{Z}[x]
\]
使得 \(|f(\alpha)|\) 很小（理想情况是 \(f(\alpha)=0\)，
即找出 \(\alpha\) 的最小多项式）。

{格的构造（How to build a lattice）}
取一组\emph{尺度} \(k\in\mathbb{N}\)（比如 \(k=6\)），并选择多项式次数上界 \(d\)。
构造如下 \((d+1)\times(d+1)\) 矩阵 \(B\) 的列向量作为格基：
\[
B=\begin{bmatrix}
1      &        &        & 0\\
       & 1      &        & 0\\
       &        & \ddots & 0\\
0      & 0      & \cdots & 1\\[2mm]
10^k   & \lfloor 10^k\alpha\rfloor & \lfloor 10^k\alpha^2\rfloor & \cdots
\end{bmatrix}
\quad(\text{上方是 }I_{d+1}, \text{ 最后一行放 }10^k\alpha^i \text{ 的取整——为了保证是整数k的大小决定了精度}).
\]
\begin{dxtips}
    其中，对角线上的元素依次代表 $x^{0}$ 到 $x^{d-1}$。  
它们的线性组合所形成的每一列，都对应于一组 $(1,\,x,\,x^{2},\,\dots,\,x^{d-1})$ 的系数。  
最后一行则是为了将这组线性组合的x直接代入 $\alpha$，使结果能够清晰直观地显示出来。我们的目的是让这组线性组合经可能的小最好是0。这里还有一个放大的好处由于最后一行远大于前面的行所以最后一行对向量的大小起决定性作用，最短的向量的最后一行必然是所有最后一行元素最小的。
\end{dxtips}


*{用 LLL 找“短向量” \(\Rightarrow\) 找多项式系数}
对基 \(B\) 运行 LLL 算法，得到一个“较短且近似正交”的等价基。
LLL 输出基中的\emph{最短向量}或若干很短的向量，
其整数系数 \(z=(c_0,\dots,c_d)\) 往往满足
\(|f(\alpha)|=|\sum c_i\alpha^i|\) 很小。
于是我们就得到了一个候选多项式
\[
f(x)=c_0+c_1x+\cdots+c_dx^d\in\mathbb{Z}[x],
\quad\text{且}\quad |f(\alpha)|\ \text{很小}.
\]
若 \(f(\alpha)\approx0\)，可尝试化简（去公因数）并检验是否确为最小多项式。

{具体例子：\(\alpha\approx\sqrt2\), \(d=2\), \(k=6\)}
取
\[
\alpha=1.414213,\quad d=2,\quad k=6.
\]
构造
\[
B=
\begin{bmatrix}
1&0&0\\
0&1&0\\
0&0&1\\
10^6&\lfloor10^6\alpha\rfloor&\lfloor10^6\alpha^2\rfloor
\end{bmatrix}
=
\begin{bmatrix}
1&0&0\\
0&1&0\\
0&0&1\\
1000000&1414213&2000000
\end{bmatrix}
\quad(\text{因为 }\alpha^2\approx2,\ \lfloor10^6\alpha^2\rfloor\approx2000000).
\]
对该 \(B\) 做 LLL 约化，得到的一个最短向量（或很短向量）可能是
\[
(-2,\,0,\,1)^\top.
\]
它对应的多项式为
\[
f(x)=(-2)+0\cdot x+1\cdot x^2=x^2-2,
\]
并且
\[
f(\alpha)=\alpha^2-2\approx0.
\]
这正是 \(\sqrt2\) 的最小多项式。

{总结}
\begin{itemize}
  \item 将“寻找整系数多项式 \(f\) 使 \(f(\alpha)\) 小”转化为“寻找格中的\emph{短向量}”；
  \item 通过构造含 \(I_{d+1}\) 与 \(\lfloor10^k\alpha^i\rfloor\) 的基矩阵 \(B\)，使每个整数系数向量 \(z\) 对应一个多项式；
  \item 用 LLL 在多项式时间内找到短向量，其系数即为候选多项式；
  \item 在 \(\alpha\approx\sqrt2\) 的例子中，LLL 返回 \((-2,0,1)\)，对应 \(x^2-2\)。
\end{itemize}

如果生成的不是最短多项式就是小数点后面精度不够。最短多项式就是，项数最少的多项式，LLL查找不同数字间的整数关系。整数本身就是一种简化。
\begin{example}[Merkle--Hellman 背包加密示例]

\textbf{1. 密钥生成}  

选择一个超递增序列：
\[
w = (2,\,3,\,7,\,14,\,30).
\]

取模数 $q=61$（大于 $\sum w_i=56$），并选取 $r=17$ 与 $q$ 互素。  

计算公钥：
\[
b_i \equiv r \cdot w_i \pmod q,
\]
得到
\[
b=(34,\,51,\,58,\,54,\,22).
\]

因此：  
- 私钥为 $(w,q,r)$；  
- 公钥为 $b=(34,51,58,54,22)$。  
\textbf{2. 加密}  

明文消息为比特串 $(1,0,1,0,1)$。  

计算密文：
\[
c = \sum_{i=1}^5 m_i b_i = 1\cdot34 + 0\cdot51 + 1\cdot58 + 0\cdot54 + 1\cdot22 = 114.
\]

因此密文为：
\[
c=114.
\]
\textbf{3. 解密}  

首先计算 $r$ 在模 $q$ 下的逆元：
\[
r^{-1} \equiv 18 \pmod{61}, \quad \text{因为 } 17\cdot18 \equiv 1 \pmod{61}.
\]

然后：
\[
c' \equiv c \cdot r^{-1} \pmod q = 114 \cdot 18 \equiv 37 \pmod{61}.
\]

在超递增序列 $w=(2,3,7,14,30)$ 中解方程
\[
\sum_{i=1}^5 m_i w_i = 37.
\]

使用贪心法：  
- $30 \leq 37 \implies m_5=1$，余数 $7$；  
- $14>7 \implies m_4=0$；  
- $7=7 \implies m_3=1$，余数 $0$；  
- 余下均为 $0$。  

因此恢复的明文为：
\[
(m_1,m_2,m_3,m_4,m_5) = (1,0,1,0,1).
\]

\end{example}

\begin{remark}
Merkle--Hellman 背包加密体系最初被认为是安全的，因为从公钥 $b=(b_1,\dots,b_n)$ 
和密文 $c=\sum m_i b_i$ 恢复比特向量 $(m_1,\dots,m_n)$ 等价于求解
\[
\sum_{i=1}^n m_i b_i = c, \qquad m_i \in \{0,1\},
\]
这正是著名的 \emph{子集和问题}（subset sum problem），在一般情况下是 NP-困难的。  

然而，LLL 格基约化算法提供了一种有效的攻击途径：  
通过构造一个 $(n+1)$ 维格，其中包含向量
\[
(a_1,\dots,a_n,\,c),
\]
LLL 可以在多项式时间内找到短向量。  
这些短向量往往正好对应于满足上述子集和关系的比特向量 $(m_1,\dots,m_n)$。  

因此，原本依赖于子集和难解性的 Merkle--Hellman 背包体系，
在 LLL 算法出现之后不再安全，很快被实际攻破。  
这也是格密码学对现代密码体系影响的重要历史事件之一。
\end{remark}
\begin{definition}[时间复杂度]
在计算机科学中，算法的\emph{时间复杂度}是指：  
算法在最坏情况或平均情况下所需的基本运算次数，作为输入规模 $n$ 的函数。  
换句话说，时间复杂度研究的是当 $n$ 越来越大时，算法运行时间如何随 $n$ 增长。

\medskip

设算法 $A$ 在输入规模为 $n$ 时，所需的基本操作次数为 $T(n)$。  
如果存在常数 $c>0$ 和 $n_0$，使得对所有 $n \geq n_0$ 都有
\[
T(n) \leq c \cdot f(n),
\]
则称该算法的时间复杂度是 $O(f(n))$。  
其中 $f(n)$ 是一个常见的函数，例如 $n,\; n^2,\; n\log n$ 等。
\end{definition}
\section{LLL的时间复杂度}
\begin{theorem}[LLL 算法的时间复杂度（经典结论）]
设输入是 $n$ 维整数格的基 $B\in\mathbb{Z}^{n\times n}$，其元素绝对值不超过 $B$。
对固定的 $\delta\in(1/4,1)$，标准 LLL 算法在位操作意义下的时间复杂度满足
\[
T(n,B,\delta)\;=\;\widetilde{O}\!\bigl(n^{4}\cdot \log B\bigr),
\]
更精细的经典上界常写作
\[
O\!\bigl(n^{4}\cdot \log^{3} B\bigr),
\]
其中 $\widetilde{O}$ 省略了多项式对数因子。实际实现常优于理论上界。
\end{theorem}
GSO就是施密特正交化
\begin{table}[h]
\centering
\begin{tabular}{|c|l|c|l|}
\hline
步骤 & 动作 & 单次代价 & 备注 \\
\hline
1 & 初始化 GSO （计算 $\mu_{i,j}, \|b_j^*\|^2$） & $O(n^3)$ & 一次性操作 \\
\hline
2 & 尺寸约化（对 $j<i$：算 $\mu_{i,j}$，更新 $b_i$） & $O(n^2)$ & 内积 $O(n)$ + 向量更新 $O(n)$，共 $i-1\le n$ 次 \\
\hline
3 & Lovász 条件检查 & $O(1)$ & 只用已存的 $\mu,\|b^*\|^2$ \\
\hline
4 & 交换并更新 GSO （若条件不满足） & $O(n^2)$ & 增量更新避免 $O(n^3)$ \\
\hline
循环次数 & 交换步数 $\le O(n^2 \log B)$ & —— & 势函数下降法保证 \\
\hline
总复杂度 & \multicolumn{3}{c|}{$O(n^4 \log^3 B)$ （其中 $\log^3 B$ 来自大整数算术位长）} \\
\hline
\end{tabular}
\caption{LLL 算法每步操作及时间复杂度}
\end{table}
\begin{theorem}[LLL 算法的复杂度上界]
设输入基为 $B \in \mathbb{Z}^{n \times n}$，其元素的绝对值不超过 $B$。  
则 Lenstra--Lenstra--Lovász (1982) 的分析表明：  
\begin{itemize}
  \item 所有中间数的位长增长是多项式有界，不会发生指数爆炸；
  \item 在保守估计下，每次算术操作的位复杂度需要额外考虑 $(\log B)^2$ 或 $(\log B)^3$。
\end{itemize}
因此，LLL 算法的总时间复杂度满足
\[
T(n,B) = O\!\left(n^4 \cdot \log^3 B\right).
\]
\end{theorem}

\begin{proposition}[对 $\log^3 B$ 的拆解理解]
复杂度中的 $\log^3 B$ 可以解释为以下三个因素的叠加：
\begin{enumerate}
  \item 第一个 $\log B$：来自输入规模，即最大数 $B$ 的比特长度；
  \item 第二个 $\log B$：来自大整数算术运算（乘法、除法）的代价；
  \item 第三个 $\log B$：来自迭代过程中中间数的膨胀，虽数值变大但仍在多项式范围内。
\end{enumerate}
因此最终得到 $\log^3 B$。
\end{proposition}
